\chapter{Measuring before modelling}
\label{ch:first}

\section{The goal}

A trained model produces a number. Whether that number means anything depends on
three things that have nothing to do with the model: which saved copy of the run
is being read, where the day was cut in two, and how much room the benchmark
leaves in the first place.

Each of the three cost me a wrong conclusion before it was understood, and each
turned out to be worth more, in correct answers, than most of the modelling ideas
that came later. This chapter states the three rules and uses the early runs as
the evidence for them.

\paragraph{Where the numbers in this chapter come from.} From the three earliest
sets of people of Table~\ref{tab:datasets}: \Draw{} taken exactly as the operator
sent it, \Dfilt{} once people with too few records were removed, and \Dmob{} once
only people who genuinely move are kept. Each uses \textbf{400 people for
training} and \textbf{100 different people}, never seen during training, for
grading. The one-line rule scores 69.2\,\%, 79.0\,\% and 45.5\,\% on them
respectively, and that difference between the three is itself the subject of the
third rule below.

\section{How a prediction is scored}

Everything in this report is scored the same way, and the procedure never
changes. It has three parts.

\textbf{First, the day is cut in two.} Each test person's day is split once into
a past, which the model may read, and a rest, which it has to produce.
Figure~\ref{fig:protocol}(a) draws the cut, and the default is 80/20: the model
sees the first four fifths of somebody's day and is graded on the last fifth.

\textbf{Second, the rest is answered one moment at a time.} The model is never
asked to invent a whole afternoon. It is asked for one antenna, then told what
really happened, then asked for the next one, as in
Figure~\ref{fig:protocol}(b). So it can neither see the future nor compound its
own mistakes. The one-line rule is run through exactly the same moments with
exactly the same history, and answers the last antenna it was shown.

% F5 -- Where a day is cut, and how a single position is scored.
\begin{figure}[!htbp]
\centering
\fitwidth{%
\begin{tikzpicture}[font=\small,
  ar/.style={-{Stealth[length=4pt]},clrule},
  ctx/.style={draw=clrule,fill=clsoft,rounded corners=2pt,inner sep=4pt,
              font=\scriptsize,text width=3.1cm,align=center},
  prd/.style={draw=clmodel,fill=clmodel!8,rounded corners=2pt,inner sep=4pt,
              font=\scriptsize,text width=2.2cm,align=center},
]
% ---- (a) the split ------------------------------------------------------
\node[anchor=west,font=\footnotesize\bfseries,text=clmodel] at (0,1.05)
  {(a) Each person's day is cut in two, once and for all};
\fill[clmodel!22] (0,0) rectangle (11.1,0.65);
\fill[clink!25]  (11.1,0) rectangle (13.9,0.65);
\draw[clrule] (0,0) rectangle (13.9,0.65);
\node[font=\footnotesize] at (5.55,0.33)
  {\textbf{THE PAST}, 80\,\% of the day, given to the model};
\node[font=\footnotesize] at (12.5,0.33) {\textbf{THE REST}, 20\,\%};
\node[font=\scriptsize,anchor=north,text=clrule] at (6.95,-0.10)
  {where this cut falls is called the \emph{context fraction}; several
   experiments move it from 30\,\% to 90\,\%};

% ---- (b) scoring one position at a time ---------------------------------
\node[anchor=west,font=\footnotesize\bfseries,text=clmodel] at (0,-1.20)
  {(b) The rest is then answered one moment at a time, never in one go};
\begin{scope}[yshift=-1.60cm]
  \node[ctx] (b0) at (2.6,-0.35)  {given events 1 to 40};
  \node[prd] (p0) at (5.9,-0.35)  {name event 41};
  \node[ctx] (b1) at (2.6,-1.10)  {given events 1 to 41};
  \node[prd] (p1) at (5.9,-1.10)  {name event 42};
  \node[ctx] (b2) at (2.6,-1.85)  {given events 1 to 42};
  \node[prd] (p2) at (5.9,-1.85)  {name event 43};
  \draw[ar] (b0) -- (p0); \draw[ar] (b1) -- (p1); \draw[ar] (b2) -- (p2);
  \node[font=\scriptsize,anchor=west,text width=6.6cm,align=left] at (7.6,-1.10)
    {Each time, the model is handed the person's \emph{real} past, not its own
     previous answers. So it can neither see the future nor go astray on its own
     earlier mistakes. The one-line rule answers exactly the same moments, and
     simply repeats the last antenna in the history it was given.};
\end{scope}
\end{tikzpicture}}
\caption{Where the day is cut, and how each moment is answered}
\label{fig:protocol}
\figtext{%
\figpara{Panel (a) is the cut; panel (b) is how each moment is answered.}{Each
test person's day is split once into a part the model may read and a part it has
to produce, the default being 80/20. Where that cut falls is part of the
\emph{measurement} rather than a property of the model, which is why it is drawn
here (Chapter~\ref{ch:first}). The model is then not asked to invent a whole
afternoon: it is asked for one antenna, told what really happened, and asked for
the next, so it can neither see the future nor compound its own errors. The
one-line rule is run through exactly the same moments with exactly the same
history and simply answers the last antenna it was shown, which is what makes the
two numbers comparable at all.}}
\end{figure}


\textbf{Third, the answers are counted.} This is where a percentage comes from,
and it is worth spelling out because every figure in the rest of this document is
one of these counts. At each graded moment, the model produces not a single
answer but a ranking of all 369 antennas, most likely first. The moment counts as
correct if the antenna the person really visited is at the top of that ranking.
The score is then the number of correct moments divided by the number of graded
moments: ten moments, seven correct, 70\,\%.

Reading the ranking one place down gives \textbf{top-3}, and further down
\textbf{top-5}. Figure~\ref{fig:topk} shows all of it on one picture.

% F5e -- What top-1, top-3 and top-5 mean, and how a score is computed from them.
\begin{figure}[!htbp]
\centering
\fitwidth{%
\begin{tikzpicture}[font=\small,
  cand/.style={draw=clrule,fill=white,rounded corners=2pt,align=center,
               inner sep=4pt,text width=2.15cm,font=\scriptsize},
  hit/.style={cand,draw=clgain,fill=clgain!14,thick},
  brk/.style={clmodel,thick,|-|},
  cell/.style={draw=clrule,minimum width=0.85cm,minimum height=0.55cm,
               inner sep=0pt,font=\scriptsize},
  ok/.style={cell,draw=clgain,fill=clgain!40},
  ko/.style={cell,draw=clnonsig,fill=white},
]
% ================= (a) one prediction, ranked =================
\node[anchor=west,font=\footnotesize\bfseries,text=clmodel] at (-0.35,3.55)
  {(a) The model never answers once. It ranks all 369 antennas.};
\node[anchor=west,font=\scriptsize,text=clrule] at (-0.35,3.18)
  {Here are its first five, for one moment of one person's day, with how sure it
   is of each.};

% the three depth markers, stacked above the row
\draw[brk] (-0.05,2.62) -- (2.35,2.62);
\node[font=\scriptsize,text=clmodel,anchor=west] at (2.55,2.62) {\textbf{top-1}};
\draw[brk] (-0.05,2.20) -- (7.45,2.20);
\node[font=\scriptsize,text=clmodel,anchor=west] at (7.65,2.20) {\textbf{top-3}};
\draw[brk] (-0.05,1.78) -- (12.55,1.78);
\node[font=\scriptsize,text=clmodel,anchor=west] at (12.75,1.78) {\textbf{top-5}};

\node[cand] (k1) at (1.15,0.95) {\#1\\ \cid{UKVBER101}\\ 52\,\%};
\node[cand] (k2) at (3.70,0.95) {\#2\\ \cid{UKVTGM104}\\ 23\,\%};
\node[hit]  (k3) at (6.25,0.95) {\#3\\ \cid{BKVPAN1}\\ 13\,\%};
\node[cand] (k4) at (8.80,0.95) {\#4\\ \cid{UKVDVO104}\\ 7\,\%};
\node[cand] (k5) at (11.35,0.95){\#5\\ \cid{BKVBOH1}\\ 5\,\%};
\node[font=\scriptsize,text=clrule,anchor=west] at (12.75,0.95)
  {\dots\ 364 more};

\node[font=\scriptsize,text=clgain,anchor=north,align=center] at (6.25,0.28)
  {$\uparrow$ this is the antenna the person really visited};

\node[font=\scriptsize,anchor=west] at (-0.35,-0.55)
  {\textcolor{clloss}{\textbf{top-1: wrong}}, it is not first. \qquad
   \textcolor{clgain}{\textbf{top-3: right}}, it is in the first three. \qquad
   \textcolor{clgain}{\textbf{top-5: right}}, it is in the first five.};

% ================= (b) how a score is computed =================
\draw[clrule,dotted] (-0.45,-1.05) -- (14.4,-1.05);
\node[anchor=west,font=\footnotesize\bfseries,text=clmodel] at (-0.35,-1.45)
  {(b) The score is then just a count, over every graded moment of every test
   person.};

\foreach \i/\st in {0/ok,1/ko,2/ok,3/ok,4/ko,5/ok,6/ok,7/ok,8/ko,9/ok}{
  \node[\st] at (0.1+\i*0.92,-2.25) {};
}
\node[font=\scriptsize,anchor=west,text width=6.2cm] at (9.4,-2.25)
  {shaded = right, empty = wrong\\[2pt]
   7 right out of 10 graded moments\\ $\Rightarrow$ a score of \textbf{70\,\%}};
\node[font=\scriptsize,text=clrule,anchor=west] at (-0.35,-2.95)
  {The one-line rule is graded on the very same moments, and its own count is
   what my score is compared against.};
\end{tikzpicture}}
\caption{Top-1, top-3 and top-5, and the way a score is obtained}
\label{fig:topk}
\figtext{%
\figpara{Panel (a): the model does not give an answer, it gives an order.}{Asked
what comes next, it assigns a degree of confidence to all 369 antennas and sorts
them, most likely first: here 52\,\% for \cid{UKVBER101}, 23\,\% for
\cid{UKVTGM104}, and so on. \textbf{Top-1} asks whether the right answer is in
first place, \textbf{top-3} whether it is anywhere in the first three,
\textbf{top-5} in the first five. In the example the right answer sits third, so
the same prediction fails at top-1 and succeeds at the other two. Reporting three
depths is the convention of this literature rather than a choice of mine
\cite{luca2021survey,qin2025mobllm}.}

\figpara{Panel (b): where a percentage comes from.}{The score is a plain count.
Every graded moment is answered, marked right or wrong at the chosen depth, and
the score is the number of right answers over the number of moments: ten moments,
seven right, 70\,\%. The one-line rule answers the very same moments and gets its
own count, and every result in this report is the difference between the two.}}

\end{figure}


The three depths are not three models; they are the same predictions read at
three levels of strictness. Reporting all three is the convention of the mobility
literature rather than an idea of mine \cite{luca2021survey,qin2025mobllm}, and it
exists because an operator that can prepare three antennas in advance is served
by a correct answer in third place. Top-1 is the strict headline used throughout
this report; the other two say whether a model that misses is missing narrowly or
wildly.

\section{Rule 1: never read the last pass of a run}

A language model fine-tuned on a few hundred short days will end up learning them
by heart.\footnote{\textbf{Overfitting}: the point past which a model stops
learning the regularity of the data and starts memorising the examples it was
shown.} That is not a hypothesis; it is what the curves do.

The clearest evidence in this project is a single run, deliberately allowed to
continue to a hard ceiling of forty-nine passes over the data. Its training score
improves the whole way through, but its score on people it had never seen climbs
for nine passes, peaks, and then declines for the remaining forty.
Figure~\ref{fig:validation-curves}(a) shows both halves of that sentence at once.

% F6 -- Four runs, four datasets, the same shape: a plateau under the baseline.
\begin{figure}[t]
\centering
\plotlink{\nbllm}{\fitwidth{%
\begin{tikzpicture}
\begin{groupplot}[
  group style={group size=2 by 2, horizontal sep=1.6cm, vertical sep=2.0cm,
               xlabels at=edge bottom, ylabels at=edge left},
  rep, width=0.46\linewidth, height=4.3cm,
  xlabel={epoch}, ylabel={validation top-1 (\%)},
  ymin=0, ymax=100, ytick={0,20,...,100},
  no markers, title style={yshift=2pt},
]
% ---------- (a) D1, unfiltered day, 500 users, 49 epochs ----------
\nextgroupplot[title={\footnotesize (a) \sraw, 49 epochs},
               xmin=0, xmax=50, xtick={1,10,20,30,40,49}]
\addplot[clbase,dashed,thick,domain=0:50,samples=2]{69.2};
\addplot[clmodel,thick,mark=*,mark size=0.8pt] coordinates {
(1,18.9)(2,53.8)(3,55.8)(4,61.5)(5,56.4)(6,63.9)(7,62.7)(8,60.3)(9,65.0)(10,61.8)
(11,60.8)(12,59.6)(13,59.8)(14,61.9)(15,58.7)(16,59.8)(17,60.6)(18,61.1)(19,60.2)(20,55.8)
(21,53.3)(22,53.5)(23,53.0)(24,50.7)(25,50.2)(26,50.5)(27,50.7)(28,47.9)(29,49.0)(30,47.8)
(31,46.6)(32,47.0)(33,48.0)(34,48.1)(35,47.2)(36,46.7)(37,48.3)(38,49.0)(39,47.6)(40,46.9)
(41,47.4)(42,46.9)(43,46.9)(44,47.9)(45,47.4)(46,47.4)(47,47.9)(48,48.1)(49,48.1)};
\node[star,star points=5,fill=clgain,draw=none,inner sep=1.6pt] at (axis cs:9,65.0) {};
\node[font=\scriptsize,anchor=south west,text=clgain] at (axis cs:9,66) {peak, epoch 9};
\node[font=\scriptsize,anchor=south east,text=clbase] at (axis cs:49,70) {baseline 69.2\,\%};
%
% ---------- (b) D2, filtered, 14 epochs ----------
\nextgroupplot[title={\footnotesize (b) \sfilt, 14 epochs},
               xmin=0, xmax=15, xtick={1,4,...,14}]
\addplot[clbase,dashed,thick,domain=0:15,samples=2]{79.0};
\addplot[clmodel,thick,mark=*,mark size=1pt] coordinates {
(1,56.1)(2,72.2)(3,73.2)(4,73.5)(5,71.8)(6,73.8)(7,71.8)(8,71.8)(9,69.8)(10,70.8)
(11,69.1)(12,67.6)(13,68.2)(14,67.2)};
\node[star,star points=5,fill=clgain,draw=none,inner sep=1.6pt] at (axis cs:6,73.8) {};
\node[font=\scriptsize,anchor=south,text=clgain] at (axis cs:6,75) {epoch 6};
\node[font=\scriptsize,anchor=south east,text=clbase] at (axis cs:14.5,80) {79.0\,\%};
%
% ---------- (c) D3, mobile users, 9 epochs ----------
\nextgroupplot[title={\footnotesize (c) \smob, 9 epochs},
               xmin=0, xmax=10, xtick={1,3,...,9}]
\addplot[clbase,dashed,thick,domain=0:10,samples=2]{43.5};
\addplot[clmodel,thick,mark=*,mark size=1pt] coordinates {
(1,19.8)(2,40.2)(3,42.3)(4,42.5)(5,43.0)(6,42.3)(7,39.9)(8,38.7)(9,38.4)};
\node[star,star points=5,fill=clgain,draw=none,inner sep=1.6pt] at (axis cs:5,43.0) {};
\node[font=\scriptsize,anchor=south,text=clgain] at (axis cs:5,44.5) {epoch 5};
\node[font=\scriptsize,anchor=north east,text=clbase] at (axis cs:9.7,41.5) {43.5\,\%};
%
% ---------- (d) D3 + hour tokens, 11 epochs ----------
\nextgroupplot[title={\footnotesize (d) \smob + hour tokens, 11 epochs},
               xmin=0, xmax=12, xtick={1,3,...,11}]
\addplot[clbase,dashed,thick,domain=0:12,samples=2]{45.5};
\addplot[clmodel,thick,mark=*,mark size=1pt] coordinates {
(1,21.1)(2,39.0)(3,41.9)(4,41.1)(5,42.2)(6,41.3)(7,40.0)(8,39.4)(9,38.1)(10,38.8)(11,37.1)};
\node[star,star points=5,fill=clgain,draw=none,inner sep=1.6pt] at (axis cs:5,42.2) {};
\node[font=\scriptsize,anchor=south,text=clgain] at (axis cs:5,43.6) {epoch 5};
\node[font=\scriptsize,anchor=north east,text=clbase] at (axis cs:11.7,43.5) {45.5\,\%};
\end{groupplot}
\end{tikzpicture}}}
\caption{Validation accuracy per epoch on four datasets}
\label{fig:validation-curves}
\figtext{%
The same shape, four times, on four different datasets. In each panel the
horizontal axis is the training epoch and the vertical axis is top-1 accuracy
on validation users the run never trains on; the solid line is the model, the
star its best epoch, the dashed line that dataset's persistence baseline.
Validation climbs for four to nine epochs, peaks, then decays for tens of
epochs while the training loss keeps falling. In panel~(a) the peak is
65.0\,\% at epoch~9 against 48.1\,\% at the last epoch, so reading the end of
a run instead of its best point costs sixteen points and forty epochs of GPU
time. Note that the dashed line moves between panels, 69.2\,\% on the
unfiltered day and 43.5\,\% once the corpus is restricted to mobile users,
because the users themselves are different; the panels are comparable only on
the distance between the two lines, and in all four that distance is negative.
Per-epoch values are recovered by digitising each run's own validation curve,
with the vertical axis calibrated on the two values printed on that curve, so
an individual point carries about $\pm1\pt$ of reading error; the figures
quoted in the text are the reported ones. Qwen2.5-0.5B + QLoRA throughout.}
\figsource{\nbllm}{train\_cellid\_llm.ipynb}{the training loop whose per-epoch validation scores these curves are}
\end{figure}


The important part is not that this happened, which was expected, but how much it
costs to ignore it. The copy saved at the end of that run scored 48.7\,\% on the
held-out people. The copy saved at the peak scored 64.2\,\%. Reading the end of
training instead of its best point therefore throws away \textbf{15.5 correct
answers per hundred}, and spends forty passes of waiting time doing so.

Acting on that, keeping the best copy, stopping on a plateau, and adding three
mild settings that make memorisation harder,\footnote{Three standard knobs.
\textbf{Dropout}: a tenth of the added numbers are set to zero at random at each
step, so the model cannot lean on any single one. \textbf{Learning rate}: how big
a step to take when adjusting a number. \textbf{Weight decay}: a mild pull of
every trainable number back towards zero.} moved the score on unseen people from
48.7\,\% to \textbf{64.4\,\%}. That is 15.7 more correct answers per hundred for
the saving rule and the three settings together. It also cut the waiting time
from seventy-two minutes to twenty-nine, and lifted top-3 from 76.2\,\% to
82.7\,\%.

Panels (b), (c) and (d) of the same figure confirm that this shape is a property
of the task rather than of one run: on three further sets of people, with rules
ranging from 43.5\,\% to 79.0\,\%, the peak arrives between the fourth and sixth
pass and then decays. The peak is always early, which is why every run in this
report is stopped early and read at its best copy. Two runs are then always
compared at their respective peaks, which is the only comparison that is not an
accident of where each one happened to stop.

\section{Rule 2: where the day is cut belongs to the ruler}

The share of the day given to the model, called the \emph{context
fraction},\footnote{\textbf{Context fraction}: the share of a person's day handed
to the model before it must produce the rest. Distinct from the \emph{context
window} of Chapter~\ref{ch:tools}, which is a hard architectural limit.} looks
like a setting of the experiment. It is not. It is part of the measurement.

Moving the cut changes which antenna counts as ``the last one seen'', so it
changes the one-line rule too. A model graded at 90\,\% context is not a better
model than the same model graded at 30\,\%; it is the same model measured against
an easier ruler.

Table~\ref{tab:context-sweep} is the evidence, and it contains the mistake this
rule exists to prevent. On \Dfilt{} the model reaches 80.1\,\% at the reference
80/20 split, against a rule that scores 79.0\,\%: \textbf{1.1 more correct
answers per hundred}, and the first time in the project that a trained model
cleared the one-line rule. Read on its own that looks like a result. Read across
the three rows it is far weaker: at 70\,\% context the model trails by 0.3, and
at 90\,\% context, where its own score is the \emph{highest} of the three, it
trails by 2.4. One split out of three passes.

% T3 -- The context/target split swept on two datasets (slides 27 and 37).
\begin{table}[!htbp]
\centering\small
\caption{Moving the cut moves the ruler with it}
\label{tab:context-sweep}
\begin{tabular}{@{}lrrrrrrc@{}}
\toprule
\textbf{past / rest} & \textbf{Valid.} & \textbf{Model} & \textbf{Rule} & \textbf{Distance} & \textbf{top-3} & \textbf{top-5} & \textbf{Beats it?} \\
\midrule
\multicolumn{8}{@{}l}{\textit{\sfilt}, three separate trainings, one per cut}\\
\addlinespace[2pt]
70 / 30          & 72.3\,\% & 76.5\,\% & 76.8\,\% & $-0.3$ & 87.1\,\% & 90.1\,\% & no \\
\textbf{80 / 20} & 73.8\,\% & \textbf{80.1\,\%} & 79.0\,\% & $\mathbf{+1.1}$ & 89.1\,\% & 91.6\,\% & \textbf{yes} \\
90 / 10          & 75.7\,\% & 77.7\,\% & 80.1\,\% & $-2.4$ & 87.2\,\% & 90.5\,\% & no \\
\midrule
\multicolumn{8}{@{}l}{\textit{\smob}, one training at 80/20, evaluated at seven cuts}\\
\addlinespace[2pt]
30 / 70 & 43.0\,\% & 41.5\,\% & 42.9\,\% & $-0.9$ & 62.1\,\% & 69.9\,\% & no \\
40 / 60 & 43.0\,\% & 42.0\,\% & 43.2\,\% & $-0.6$ & 62.1\,\% & 69.7\,\% & no \\
50 / 50 & 43.0\,\% & 42.0\,\% & 43.3\,\% & $-0.7$ & 62.0\,\% & 69.4\,\% & no \\
60 / 40 & 43.0\,\% & 42.7\,\% & 44.0\,\% & $-0.9$ & 62.1\,\% & 69.5\,\% & no \\
70 / 30 & 43.0\,\% & 43.6\,\% & 44.3\,\% & $\mathbf{-0.3}$ & 62.1\,\% & 69.5\,\% & no \\
80 / 20 & 43.0\,\% & 44.6\,\% & 45.5\,\% & $-0.9$ & 63.9\,\% & 70.7\,\% & no \\
90 / 10 & 43.0\,\% & 47.6\,\% & 50.9\,\% & $-2.1$ & 68.0\,\% & 74.3\,\% & no \\
\bottomrule
\end{tabular}

\figtext{%
\figpara{Table~\ref{tab:context-sweep}: parameters and column descriptions.}{Only where each person's day
is cut into the part the model reads and the part it must produce: ``70 / 30''
means the model sees the first 70\,\% and is graded on the last 30\,\% in Table~\ref{tab:context-sweep}.
\textbf{Valid.} is the score on people held back to decide when to stop training,
\textbf{Model} the score on the test people, \textbf{Rule} what ``repeat the last
antenna'' scores on those very same moments, and \textbf{Distance} the model minus
the rule, which is the only column that means anything.}

\figpara{Table~\ref{tab:context-sweep}: persistence baseline variation across context fractions.}{The \textbf{Rule} column in Table~\ref{tab:context-sweep} changes from row
to row, because moving the cut changes which antenna counts as ``the last one
seen''. On \sfilt{} exactly one cut out of three clears the rule, by 1.1, while
the other two lose ground by 0.3 and 2.4; one favourable cell out of three is not
a result, and the next iteration made it disappear. On \smob{} no cut clears the
rule at all: the model climbs steadily from 41.5 to 47.6\,\% as it is given more
history, which looks like progress until one notices the rule climbing just as
fast. Each \sfilt{} row is a separate training run; the \smob{} block is one model
read seven times, and the giveaway is its identical \textbf{Valid.} column.}}

\figsource{\nbllm}{train\_cellid\_llm.ipynb}{the context-fraction sweeps}
\end{table}


The lower block of the same table settles it. On \Dmob{} a single trained model
was graded at seven cuts from 30\,\% to 90\,\%. Its score rises steadily from
41.5\,\% to 47.6\,\%, and the one-line rule rises with it, in lockstep, at every
single value. No cut beats the rule; the closest is 0.3 short and the worst,
again at 90\,\% context, is 2.1 short. More history helps the model and helps
``repeat the last antenna'' by the same amount, so a rising accuracy curve across
cuts carries no information about the model at all.

One clarification belongs here, because the two blocks of that table are not the
same experiment. The \sfilt{} rows are three separate trainings, one per cut. The
\smob{} rows are one training at 80/20, graded seven times, so what they measure
is how well a model tolerates being handed more or less history than it was
trained on, not which cut trains the best model. Both questions are legitimate;
they are not interchangeable, and the second is the one this report can answer.

\section{Rule 3: a benchmark has to leave room}

The third rule is the one that took longest to accept.

On a day file taken as it comes, ``repeat the last antenna'' is right 69.2\,\% of
the time. There is almost nothing above that for a model to win, and most of what
is left is not mobility at all: it is people with three records from a single
antenna. A model scoring 64.4\,\% on such a file has not failed at prediction. It
has been measured on a benchmark with no room in it.

The same trained architecture, on three sets of people, faces rules of 69.2\,\%,
79.0\,\% and 45.5\,\%. Nothing about the rule changed between those three
numbers; the people did. That is why the data decision of Chapter~\ref{ch:data}
was worth more than any setting: it is what created a benchmark on which the
question could be measured at all.

\section{What the three rules settled}

Taken together they closed the ``tune it harder'' avenue. Once runs were read at
their peak, once the cut was recognised as part of the ruler, and once the people
had been chosen so that the yardstick was beatable in principle, the model still
sat within a few correct answers of a one-line rule, usually below it, at every
cut and on every set of people tried.

The example that decided what to do next is small enough to quote in full. On one
run, person 775719, whose past is the same antenna repeated ten times, is
predicted correctly with 100\,\% confidence. Person 3268099, whose past ends on
one antenna and whose next event is a different one entirely, is predicted as
``the same antenna again'' with 68\,\% confidence. The model was right for free
on the first and wrong for free on the second, and in neither case did training
have much to do with the outcome.

A companion experiment made the same point statistically: on a 100-person version
of the data, the score per person ranged from 10\,\% to 100\,\% and the model
never passed 42.1\,\%. One hundred people is not a population; it is a handful of
anecdotes with an average printed underneath. That is what
Chapter~\ref{ch:user} sets out to fix.

\section{Looking back}
\label{sec:retro-first}

Two things stand out with hindsight.

The first is that 1.1. I presented it as a result, and it was not one. Of the
three cuts tried, one passes and the two others lose ground, and every figure in
that table comes from a single training. One favourable cell in a table of three
is too weak to carry a conclusion, and the following iteration duly made it
disappear: the same architecture on cleaner data beat nothing at all. What I
would ask of that number now is the one thing I did not ask of it, which is to
repeat it. A second training, or the same cut on a second set of people, would
have cost half an hour and would have said immediately whether there was anything
there.

The second is that the three rules of this chapter arrived one at a time, and
each of them arrived \emph{after} a conclusion that turned out to be wrong. They
are not deep ideas; they are hygiene. A short checklist written before the first
training, read the peak and not the last pass, recompute the rule every time the
cut moves, check that the benchmark leaves room before spending a week inside it,
would have caught all three at once, and the modelling work would have started
two iterations earlier. That checklist is, in effect, what this chapter has
become.
